\documentclass[10pt,a4paper]{amsart}
\usepackage[margin=19mm]{geometry}
\usepackage{amsmath,amssymb,mathtools}
\usepackage[scr=rsfs,cal=euler]{mathalfa}
\usepackage{xfrac}
\usepackage[unicode,hidelinks,bookmarks=false]{hyperref}
\newcommand{\ie}{i.e.\ }
\newcommand{\cf}{cf.\ }
\newcommand{\resp}{respectively\ }
\newcommand{\Subsection}{\subsection}
\providecommand{\nobreakdash}{\nobreak\hbox{-}}
\setlength{\emergencystretch}{2em}
\multlinegap0pt
\usepackage{bm}
\usepackage{amssymb}
\usepackage{xcolor}
\usepackage{enumerate}
\usepackage{tikz}
\usepackage{algorithm}\usepackage{algpseudocode}
\usepackage{float}
\usepackage{dsfont}
\newtheorem{assumption}{Assumption}
\newtheorem{lemma}{Lemma}
\newcommand{\mcal}{\mathcal}
\title{Optimal Demixing of Nonparametric Densities}
\author{Jianqing Fan, Zheng Tracy Ke, Zhaoyang Shi}
\date{Source: arXiv:2603.27457v1 (CC BY 4.0)}
\begin{document}
\maketitle

\noindent\emph{Source and licence.} Optimal Demixing of Nonparametric Densities. Version arXiv:2603.27457v1 (CC BY 4.0), distributed by arXiv under the Creative Commons Attribution 4.0 International licence, http://creativecommons.org/licenses/by/4.0/, as recorded in the arXiv licence metadata for this exact version and shown on the abstract page https://arxiv.org/abs/2603.27457v1. Full licence notice: \texttt{licenses/arxiv-2603.27457-CC-BY-4.0.txt}.\par\smallskip

\noindent\textit{Editorial scope.} Counted unit: the article's lemma lem:multivariate-lemma4.3 (source0.tex lines 3882--3902). The article's complete original proof of it is delivered with it (source0.tex lines 3904--3988, one \texttt{proof} environment of 85 lines). No mathematics is written, repaired or re-derived: the statement and the proof are the article's own lines, in the article's own notation, and no retained line is substituted or re-flowed.  Retained uncounted as the source-local context the proof needs: lines 972--977; lines 1005--1014; lines 1659--1666.  Nothing was omitted from any retained span, so the package carries no omission record.  No retained line contains a citation, so the wrapper carries no bibliography and no bibliography entry.  No retained line contains a figure environment, an image-inclusion command or drawing code, so no image is delivered and the package declares no asset dependency.  Editorial formatting only, in addition: an AMS \texttt{amsart} preamble that restates the article's own declarations and macros used by the retained lines, namely the environments \texttt{assumption}, \texttt{lemma} on separate counters, together with the standard packages those lines need.  The retained environments print in this document as Assumption 1, Lemma 1 in order of appearance, whereas the article prints the same items with the numbering of its own full document.  The complete native source of the article as distributed in the arXiv e-print for this exact version is bundled unchanged as \texttt{sources/197350/source0.tex}.
\begin{assumption}[Nikol'ski  Smoothness]\label{assump1}
Fix $\beta > 0$ and $L_0>0$. Let $\lfloor\beta\rfloor$ denote the largest integer that is strictly smaller than $\beta$. For each $1\leq k\leq K$, the  $\lfloor\beta\rfloor$th order derivative of $g_k$ exists and satisfies:
  \[
  \left( \int \left| g_k^{(\lfloor\beta\rfloor)}(x + t) - g_k^{(\lfloor\beta\rfloor)}(x) \right|^2 dx \right)^{1/2} \le L_0 |t|^{\beta - \lfloor\beta\rfloor},\qquad\mbox{for all $t\in\mathbb{R}$}.
  \]
\end{assumption}

\begin{assumption} \label{assump2}
Let $\bar{g}(x)=\frac{1}{K}\sum_{k=1}^K g_k(x)dx$ and $\Sigma_{\boldsymbol g} = \frac{1}{K}\int {\boldsymbol g}(x) {\boldsymbol g}(x)' dx$. There exists a constant $c_1\in (0,1)$ such that 
\[
c_1\leq \lambda_{\min}(\Sigma_{\boldsymbol g}) \leq \lambda_{\max}(\Sigma_{\boldsymbol g})\leq c_1^{-1}, \qquad\mbox{and}\qquad \max_x \bar{g}(x) \leq c_1^{-1}. 
\]
Let $\Sigma_\Pi = (K/n)\Pi'\Pi$ and $\eta = \Pi'{\bf 1}_n$. Denote by $\eta_{\max}$ and $\eta_{\min}$ the maximum and minimum entries of $\eta$. There exists a constant $c_2\in (0,1)$ such that 
\[
c_2 <\lambda_{\min}(\Sigma_\Pi)\leq \lambda_{\max}(\Sigma_\Pi)\leq c_2^{-1}, \qquad\mbox{and}\qquad c_2K^{-1}n\leq \eta_{\min}\leq \eta_{\max}\leq c^{-1}_2K^{-1}n. 
\]
\end{assumption}

\begin{lemma} \label{lem:error1}
Under Assumptions~\ref{assump1}-\ref{assump2}, there exists a constant $C_1>0$ such that
\[
\int \|{\boldsymbol g}^*(x)-{\boldsymbol g}(x)\|^2dx\leq C_1 Kh^{2\beta}.
\]


\end{lemma}

\begin{lemma}[Multivariate analogue of Lemma \ref{lem:error1}]
\label{lem:multivariate-lemma4.3}
Let $d>1$.  For each $1\le k\le K$ let $g_k:\mathbb{R}^d\to\mathbb{R}$ be a probability density belonging to an \emph{anisotropic Nikol'ski class} with smoothness vector $\boldsymbol{\beta}=(\beta_1,\dots,\beta_d)$ and constant $L>0$ in the sense that for each multi-index $r=(r_1,\dots,r_d)$ with $0\le r_j\le\lfloor\beta_j\rfloor$ the mixed partial $\partial^r g_k$ exists in $L^2(\mathbb{R}^d)$ and for all $t=(t_1,\dots,t_d)\in\mathbb{R}^d$
\[
\Big\| \partial^r g_k(\cdot + t) - \partial^r g_k(\cdot)\Big\|_{L^2(\mathbb{R}^d)}
\le L \sum_{j=1}^d |t_j|^{\beta_j - r_j} .
\]
Recall that the multivariate kernel is the product kernel
\[
\mcal{K}_{\mathbf{h}}(u) = \prod_{j=1}^d \frac{1}{h_j} \mcal{K}_j\!\left(\frac{u_j}{h_j}\right),\qquad \mathbf{h}=(h_1,\dots,h_d),
\]
where the univariate kernel $\mcal{K}_j$ has order $\ell_j=\lfloor\beta_j\rfloor$ in each coordinate (i.e.\ $\int u^m K(u)\,du=0$ for $1\le m\le \ell_j$) and satisfies the integrability condition $\int |u|^{\beta_j}|K(u)|\,du<\infty$ for each $j$. Define the kernel-smoothed target:
\[
g_k^*(x) \;=\; \int K_{\mathbf{h}}(u)\, g_k(x+u)\,du .
\]
Then, there exists a constant $C>0$ such that
\[
\sum_{k=1}^K \int_{\mathbb{R}^d} \big( g_k^*(x) - g_k(x) \big)^2 \,dx
\;\le\; C \, K \sum_{j=1}^d h_j^{2\beta_j}.
\]
\end{lemma}

\begin{proof}
The proof follows the same high-level steps as the univariate case (Taylor expansion with cancellation of lower order terms by kernel moments and bounding the remainder by the anisotropic Nikol'ski condition), but we keep track of each coordinate's smoothness.

Fix $k$. Write the bias
\[
\Delta_k(x):=g_k^*(x)-g_k(x)=\int \mcal{K}_{\mathbf{h}}(u)\big(g_k(x+u)-g_k(x)\big)\,du.
\]
Let $\ell_j=\lfloor\beta_j\rfloor$ for $j=1,\dots,d$, and let $\ell=(\ell_1,\dots,\ell_d)$. We perform a multivariate Taylor expansion with integral-form remainder. For any multi-index $\alpha=(\alpha_1,\dots,\alpha_d)$ with $0\le\alpha_j\le\ell_j$, denote $|\alpha|=\sum_j\alpha_j$ and $\partial^\alpha$ the mixed partial derivative. The multivariate Taylor expansion about $x$ up to order $\ell$ reads
\[
g_k(x+u)-g_k(x)
= \sum_{0<|\alpha|\le|\ell|}
\frac{u^\alpha}{\alpha!}\,\partial^\alpha g_k(x)
\;+\;
R_\ell(x,u),
\]
where $u^\alpha=\prod_{j=1}^d u_j^{\alpha_j}$ and the remainder has the integral representation
\[
R_\ell(x,u)
= \sum_{|\alpha|=|\ell|}
\frac{|\ell|}{\alpha!}\int_0^1 (1-\tau)^{|\ell|-1}
\Big(\partial^\alpha g_k(x+\tau u)-\partial^\alpha g_k(x)\Big)\,d\tau\; u^\alpha .
\]
Because the kernel is a product kernel and each univariate kernel $\mcal{K}_j$ has vanishing moments up to order $\ell_j$ in coordinate \(j\), every polynomial term in the finite sum above integrates to zero against $\mcal{K}_{\mathbf{h}}(u)$. Concretely, for any multi-index $\alpha$ with $0<\alpha_j\le\ell_j$ for some $j$, by the product structure
\[
\int \mcal{K}_{\mathbf{h}}(u)\, u^\alpha\,du
=\prod_{j=1}^d \Big(\int \frac{1}{h_j} \mcal{K}_j\!\Big(\frac{u_j}{h_j}\Big) u_j^{\alpha_j}\,du_j\Big)
= \prod_{j=1}^d h_j^{\alpha_j}\Big(\int v^{\alpha_j}\mcal{K}_j(v)\,dv\Big)=0,
\]
because at least one coordinate integral vanishes. Hence all polynomial terms vanish and only the remainder contributes:
\[
\Delta_k(x)=\int \mcal{K}_{\mathbf{h}}(u)\, R_\ell(x,u)\,du.
\]

We now bound the $L^2(\mathbb{R}^d)$ norm of $\Delta_k$. Using Minkowski's integral inequality (twice) and Fubini,
\begin{align*}
\|\Delta_k\|_{L^2} 
&= \Big\| \int \mcal{K}_{\mathbf{h}}(u)\, R_\ell(\cdot,u)\,du \Big\|_{L^2}
\le \int |\mcal{K}_{\mathbf{h}}(u)|\, \|R_\ell(\cdot,u)\|_{L^2}\,du.
\end{align*}
Plug in the integral expression for $R_\ell$ and use triangle inequality and Jensen inequality (integrating in $\tau$):
\[
\|R_\ell(\cdot,u)\|_{L^2}
\le C_\ell \sum_{|\alpha|=|\ell|} |u^\alpha| \int_0^1 (1-\tau)^{|\ell|-1}
\|\partial^\alpha g_k(\cdot+\tau u)-\partial^\alpha g_k(\cdot)\|_{L^2}\,d\tau,
\]
where \(C_\ell\) is a combinatorial constant depending on \(\ell\). By the anisotropic Nikol'ski condition we have, for each such \(\alpha\),
\[
\|\partial^\alpha g_k(\cdot+\tau u)-\partial^\alpha g_k(\cdot)\|_{L^2}
\le L \sum_{j=1}^d |\tau u_j|^{\beta_j-\alpha_j}
\le L \sum_{j=1}^d |u_j|^{\beta_j-\alpha_j}.
\]
Hence, there exists a constant $C>0$ such that
\[
\|R_\ell(\cdot,u)\|_{L^2}
\le C \sum_{|\alpha|=|\ell|} |u^\alpha|\sum_{j=1}^d |u_j|^{\beta_j-\alpha_j}
\le C \sum_{j=1}^d |u_j|^{\beta_j} \cdot P(u),
\]
where \(P(u)\) is a polynomial factor in \(|u_1|,\dots,|u_d|\) whose degree depends only on \(\ell\) and \(d\). Concretely one can bound \(P(u)\le C'(1+\sum_{j}|u_j|^{\ell_j})^{|\ell|-1}\), which is harmless because of the integrability of the kernel moments assumed below.

Now, return to the $\|\Delta_k\|_{L^2}$ bound:
\[
\|\Delta_k\|_{L^2}
\le C \int |\mcal{K}_{\mathbf{h}}(u)| \sum_{j=1}^d |u_j|^{\beta_j} P(u)\,du
\le C \sum_{j=1}^d \int \prod_{m=1}^d \frac{1}{h_m}\big|\mcal{K}_j\big(\tfrac{u_m}{h_m}\big)\big|
\, |u_j|^{\beta_j} P(u)\,du.
\]
Change variables $v_m=u_m/h_m$. Then
\[
\|\Delta_k\|_{L^2}
\le C \sum_{j=1}^d h_j^{\beta_j} \int \prod_{m=1}^d |\mcal{K}_j(v_m)|\, |v_j|^{\beta_j} \tilde P(v)\,dv
\le C''\sum_{j=1}^d h_j^{\beta_j},
\]
where $\tilde P(v)$ is the rescaled polynomial and the final integral is finite by the kernel integrability assumptions $\int |v|^{\beta_j}|\mcal{K}_j(v)|\,dv<\infty$ and bounded moments. Thus
\[
\|\Delta_k\|_{L^2} \le C\sum_{j=1}^d h_j^{\beta_j}.
\]

Squaring and summing over $k=1,\dots,K$ gives
\[
\sum_{k=1}^K \|\Delta_k\|_{L^2}^2
\le K \cdot C^2 \Big(\sum_{j=1}^d h_j^{\beta_j}\Big)^2
\le C' K \sum_{j=1}^d h_j^{2\beta_j},
\]
where the last inequality follows since $(\sum_j a_j)^2 \le d\sum_j a_j^2$ and constants are absorbed into \(C'\). This is the desired bound.
\end{proof}

\end{document}
